% questions.tex —— 1977 年·北京卷·文科 · 题目内容片段（无答案）
% 由 tools/generate_projects.py 依来源总账自 1977/1977_数学真题_北京文_试题.md 生成，勿手改（改源 MD 后重跑）


\begin{gktip}
\textbf{来源说明}：本题面依据固定源 PDF《1977 北京文科数学试题》第 1 页逐题转录。源文件为 A3 单页扫描，原卷未附答案页；配套答案与解析为本库编者独立推导，\textbf{非官方标准答案}。
\end{gktip}

\par\medskip

\begin{question}
\gkti{1}计算：

\nopagebreak[4]
\[
3^0+3^{-1}-\left(1\frac79\right)^{\frac12}.
\]

\end{question}
\begin{question}
\gkti{2}化简：

\nopagebreak[4]
\[
\frac{\sqrt6+\sqrt2}{\sqrt6-\sqrt2}.
\]

\end{question}
\begin{question}
\gkti{3}解方程：

\nopagebreak[4]
\[
\frac1{x-1}+1=\frac{4x-2}{x^2-1}.
\]

\end{question}
\begin{question}
\gkti{4}不查表求 $\sin105^\circ$ 的值。

\end{question}
\begin{question}
\gkti{5}一个正三棱柱形的零件，它的高是 $10\ \mathrm{cm}$，底面边长是 $2\ \mathrm{cm}$，求它的体积。

\end{question}
\begin{question}
\gkti{6}一条直线过点 $(1,-3)$，并且与直线 $2x+y-5=0$ 平行，求这条直线的方程。

\end{question}
\begin{question}
\gkti{7}证明：等腰三角形两腰上的高相等。

\end{question}
\begin{question}
\gkti{8}为了测湖岸边 $A\text{、}B$ 两点的距离，选择一点 $C$，测得 $CA=50$ 米，$CB=30$ 米，$\angle ACB=120^\circ$，求 $AB$。

\end{question}
\begin{question}
\gkti{9}在 $2$ 和 $30$ 中间插入两个正数，这两个正数插入后使前三个数成等比数列，后三个数成等差数列，求插入的两个数。

\end{question}
\begin{question}
\gkti{10}已知二次函数 $y=x^2-4x+3$。

（1） 求出它的图象的顶点坐标和对称轴方程；

（2） 画出它的图象；

（3） 求出它的图象与直线 $y=x-3$ 的交点坐标。

\end{question}